\documentclass[10pt,a4paper]{amsart}
\usepackage[margin=19mm]{geometry}
\usepackage{amsmath,amssymb,mathtools}
\usepackage[scr=rsfs,cal=euler]{mathalfa}
\usepackage{xfrac}
\usepackage[unicode,hidelinks,bookmarks=false]{hyperref}
\newcommand{\ie}{i.e.\ }
\newcommand{\cf}{cf.\ }
\newcommand{\resp}{respectively\ }
\newcommand{\Subsection}{\subsection}
\providecommand{\nobreakdash}{\nobreak\hbox{-}}
\setlength{\emergencystretch}{2em}
\multlinegap0pt
\newtheorem{cdrthm}{Theorem}[section]
\theoremstyle{plain}
\newtheorem{thm}[cdrthm]{Theorem}
\newtheorem{lem}[cdrthm]{Lemma}
\newtheorem{prop}[cdrthm]{Proposition}
\newtheorem{cor}[cdrthm]{Corollary}
\theoremstyle{definition}
\newtheorem{defi}[cdrthm]{Definition}
\theoremstyle{remark}
\newtheorem{rmk}[cdrthm]{Remark}
\newtheorem{ex}[cdrthm]{Example}
\newtheorem{quest}[cdrthm]{Question}
\newtheorem{claim}[cdrthm]{Claim}
\title[Quasipolynomial cumulative growth for periodic graphs]{Quasipolynomial cumulative growth for strongly connected periodic graphs with C1 plus C2 less than one: original Theorem 3.4(1)(2)}
\author{Takuya Inoue and Yusuke Nakamura}
\date{Source: Algebraic Combinatorics 7 (2024), 969--1010; DOI 10.5802/alco.367; publisher version}
\begin{document}
\maketitle
\noindent\textit{Editorial scope.} The counted unit is original Theorem 3.4(1)(2) with the author's complete proof. Assertion (3) is displayed literally because it belongs to the printed statement, but it is not claimed. Uncounted context is the original Section 2 notation together with Appendix A; the uncounted required core is Appendix B, the shifted-dilation Ehrhart theorem B.4 with its complete published proof. The strongly-connected hypothesis and the exact deviation condition C1+C2<1 are unchanged, every retained passage is native author text, and printed slips are kept literal without mathematical repair.
\setcounter{section}{1}
\section{Notation and preliminaries (original \S2; uncounted context)}
\setcounter{subsection}{1}
\subsection{Graphs and walks}\label{subsection:GandW}

In this paper, a \textit{graph} means a directed weighted graph which may have loops and multiple edges. 
A graph $\Gamma = (V_{\Gamma}, E_{\Gamma}, s_{\Gamma}, t_{\Gamma}, w_{\Gamma})$ 
consists of the set $V_{\Gamma}$ of vertices, the set $E_{\Gamma}$ of edges, 
the source function $s_{\Gamma}: E_{\Gamma} \to V_{\Gamma}$, 
the target function $t_{\Gamma}: E_{\Gamma} \to V_{\Gamma}$, and 
the (integer) weight function $w_{\Gamma}: E_{\Gamma} \to \mathbb{Z}_{> 0}$. 
We often abbreviate $s_{\Gamma}$, $t_{\Gamma}$ and $w_{\Gamma}$ to $s$, $t$ and $w$ 
when no confusion can arise. 


\begin{defi}
Let $\Gamma = (V_{\Gamma}, E_{\Gamma}, s, t, w)$ be a graph. 

\begin{enumerate}
\item 
$\Gamma$ is called to be \textit{unweighted} if $w(e) = 1$ holds for every $e \in E_{\Gamma}$. 
$\Gamma$ is called to be \textit{undirected} when 
there exists an involution $E_{\Gamma} \to E_{\Gamma};\ e \mapsto e'$
such that $s(e)=t(e')$, $t(e)=s(e')$ and $w(e)=w(e')$. 
$\Gamma$ is called to be \textit{locally finite} when for all $x \in V_{\Gamma}$, 
there are only finitely many edges $e$ satisfying $s(e)=x$ and only finitely many edges $e$ satisfying $t(e)=x$. 

\item
A \textit{walk} $p$ in $\Gamma$ is a sequence $e_1 e_2 \cdots e_{\ell}$ of edges $e_i$ of $\Gamma$ satisfying 
$t(e_i) = s(e_{i+1})$ for each $i = 1, \ldots , \ell -1$. 
We define 
\[
s(p) \coloneqq  s(e_1), \quad
t(p) \coloneqq  t(e_{\ell}), \quad
w(p) \coloneqq  \sum _{i = 1} ^{\ell} w(e_i), \quad 
\operatorname{length}(p) \coloneqq  \ell. 
\]
Note that we have $w(p) = \operatorname{length}(p)$ if $\Gamma$ is unweighted. 


We say that ``$p$ is a walk from $x$ to $y$'' when $x = s(p)$ and $y = t(p)$. 
We also define the \textit{support} $\operatorname{supp}(p) \subset V_{\Gamma}$ of $p$ by
\[
\operatorname{supp}(p) \coloneqq  \{ s(e_1), t(e_1), t(e_2), \ldots , t(e_{\ell}) \} \subset V_{\Gamma}. 
\]

By convention, each vertex $v \in V_{\Gamma}$ is also considered as a walk of length $0$. 
This is called the \textit{trivial walk} at $v$ and denoted by $\emptyset _v$: i.e., 
we define 
\[
s(\emptyset _v) \coloneqq  v, \ \ 
t(\emptyset _v) \coloneqq  v, \ \ 
w(\emptyset _v) \coloneqq  0, \ \  
\operatorname{length}(\emptyset _v) \coloneqq  0, \ \ 
\operatorname{supp}(\emptyset _v) \coloneqq  \{ v \}.
\]

\item 
A \textit{path} in $\Gamma$ is a walk $e_1 \cdots e_{\ell}$ such that $s(e_1), t(e_1), t(e_2), \ldots, t(e_{\ell})$ are distinct. 
A walk of length $0$ is considered as a path. 

\item
A \textit{cycle} in $\Gamma$ is a walk $e_1 \cdots e_{\ell}$ with $s(e_1) = t(e_{\ell})$ 
such that $t(e_1), t(e_2), \ldots, t(e_{\ell})$ are distinct. 
A walk of length $0$ is not considered as a cycle. 
$\operatorname{Cyc}_{\Gamma}$ denotes the set of cycles in $\Gamma$. 

\item 
For $x, y \in V_{\Gamma}$, $d_{\Gamma}(x, y) \in \mathbb{Z}_{\ge 0} \cup \{ \infty \}$ denotes the smallest weight $w(p)$ of any walk $p$ from $x$ to $y$. 
By convention, we define $d_{\Gamma}(x, y) = \infty$ when there is no walk from $x$ to $y$. 
A graph $\Gamma$ is said to be \textit{strongly connected} 
when we have $d_{\Gamma}(x, y) < \infty$ for all $x, y \in V_{\Gamma}$. 
When $\Gamma$ is undirected, we have $d_{\Gamma}(x,y) = d_{\Gamma}(y,x)$ for all $x,y \in V_{\Gamma}$. 

\item 
$C_1(\Gamma, \mathbb{Z})$ denotes the group of $1$-chains on $\Gamma$ with coefficients in $\mathbb{Z}$, i.e., 
$C_1(\Gamma, \mathbb{Z})$ is the free abelian group generated by $E_{\Gamma}$. 
For a walk $p = e_1 \cdots e_{\ell}$ in $\Gamma$, let $\langle p \rangle$ denote the $1$-chain 
$\sum _{i = 1} ^{\ell} e_i \in C_1(\Gamma, \mathbb{Z})$. 
$H_1(\Gamma, \mathbb{Z}) \subset C_1(\Gamma, \mathbb{Z})$ denotes the $1$-st homology group, i.e., 
$H_1(\Gamma, \mathbb{Z})$ is a subgroup generated by $\langle p \rangle$ for $p \in \operatorname{Cyc}_{\Gamma}$. 
We refer the reader to~\cite{Sunada} for more detail. 
\end{enumerate}
\end{defi}

\setcounter{subsection}{2}
\subsection{Periodic graphs}

\begin{defi}\label{defi:pg}
Let $n$ be a positive integer. 
An \textit{$n$-dimensional periodic graph} $(\Gamma, L)$ is a graph $\Gamma$ and 
a free abelian group $L \simeq \mathbb{Z}^n$ of rank $n$ with the following two conditions:
\begin{itemize}
\item $L$ freely acts on both $V_\Gamma$ and $E_\Gamma$, 
and their quotients $V_\Gamma / L$ and $E_\Gamma / L$ are finite sets. 
\item This action preserves the edge relations, i.e.,\ for any $u \in L$ and $e \in E_\Gamma$, 
we have $s_\Gamma(u(e))=u(s_\Gamma(e))$, $t_\Gamma(u(e))=u(t_\Gamma(e))$ and 
$w_{\Gamma}(u(e))=w_{\Gamma}(e)$.
\end{itemize}
Then, $L$ is called the \textit{period lattice} of $\Gamma$. 
Note that $\Gamma$ automatically becomes a locally finite graph. 
\end{defi}

If $(\Gamma, L)$ is an $n$-dimensional periodic graph, 
then the \textit{quotient graph} $\Gamma /L = (V_{\Gamma/L}, E_{\Gamma/L}, s_{\Gamma/L},t_{\Gamma/L},w_{\Gamma/L})$ is defined by 
$V_{\Gamma/L} \coloneqq  V_{\Gamma} /L$, $E_{\Gamma/L} \coloneqq  E_{\Gamma} /L$, 
and the functions $s_{\Gamma/L}: E_\Gamma/L \to V_\Gamma/L$, 
$t_{\Gamma/L}: E_\Gamma/L \to V_\Gamma/L$, and 
$w_{\Gamma/L}: E_\Gamma/L \to \mathbb{Z}_{>0}$ induced from $s_{\Gamma}$, $t_{\Gamma}$, and $w_{\Gamma}$. 
Note that the functions $s_{\Gamma/L}$, $t_{\Gamma/L}$ and $w_{\Gamma/L}$ 
are well-defined due to the second condition in Definition~\ref{defi:pg}.

\begin{defi}\label{defi:peri}
Let $(\Gamma,L)$ be an $n$-dimensional periodic graph.
\begin{enumerate}
\item Since $L$ is an abelian group, we use the additive notation: 
for $u \in L$, $x \in V_{\Gamma}$, $e \in E_{\Gamma}$ and a walk $p = e_1 \cdots e_{\ell}$,
$u + x$, $u + e$ and $u + p$ denote their translations by $u$. 

\item 
For any $x \in V_{\Gamma}$ and $e \in E_{\Gamma}$, 
let $\overline{x} \in V_{\Gamma /L}$ and $\overline{e} \in E_{\Gamma /L}$ denote their images in 
$V_{\Gamma /L} = V_{\Gamma} /L$ and $E_{\Gamma /L} = E_{\Gamma} /L$. 
For a walk $p = e_1 \cdots e_{\ell}$ in $\Gamma$, 
let $\overline{p} \coloneqq  \overline{e_1} \cdots \overline{e_{\ell}}$ denote its image in $\Gamma /L$. 

\item 
When $x, y \in V_{\Gamma}$ satisfy $\overline{x} = \overline{y}$, there exists an element $u \in L$ such that $u + x = y$. 
Since the action is free, such $u \in L$ uniquely exists and is denoted by $y - x$. 

\item 
For a walk $p$ in $\Gamma$ with $\overline{s(p)} = \overline{t(p)}$, we define 
\[
\operatorname{vec}(p) \coloneqq  t(p) - s(p) \in L. 
\]
\end{enumerate}
\end{defi}

\begin{defi}\label{defi:pr}
Let $(\Gamma,L)$ be an $n$-dimensional periodic graph.
We define $L_{\mathbb{R}} \coloneqq  L \otimes _{\mathbb{Z}} \mathbb{R}$. 
\begin{enumerate}
\item 
A \textit{periodic realization} $\Phi\colon V_{\Gamma} \to L_{\mathbb{R}}$ is a map satisfying
$\Phi(u+x) = u + \Phi(x)$ for any $u \in L$ and $x \in V _{\Gamma}$. 
When we fix an injective periodic realization of $\Phi\colon V _{\Gamma} \to L_{\mathbb{R}}$, 
we sometimes identify $V _{\Gamma}$ with the subset of $L_{\mathbb{R}}$. 

\item 
Let $\Phi$ be a periodic realization of $(\Gamma, L)$. 
For an edge $e$ and a walk $p$ in $\Gamma$, we define 
\[
\operatorname{vec}_{\Phi}(e) \coloneqq  \Phi(t(e)) - \Phi(s(e)) \in L_{\mathbb{R}}, \quad
\operatorname{vec}_{\Phi}(p) \coloneqq  \Phi(t(p)) - \Phi(s(p)) \in L_{\mathbb{R}}. 
\]
It is easy to see that the value $\operatorname{vec}_{\Phi}(e) \in L_{\mathbb{R}}$ depends only on the class $\overline{e} \in E_{\Gamma /L}$, 
and therefore, 
the map 
\[
\mu _{\Phi} : E_{\Gamma /L} \to L_{\mathbb{R}};\quad \overline{e} \mapsto \operatorname{vec}_{\Phi}(e)
\]
is well-defined. 
It can be extended to a homomorphism 
\[
\mu _{\Phi}: C_1(\Gamma/L, \mathbb{Z}) \to L_{\mathbb{R}}; \quad 
\sum a_i \overline{e_i} \mapsto \sum a_i \mu _{\Phi} (\overline{e_i}). 
\]
By construction, it satisfies $\mu _{\Phi} (\langle \overline{p} \rangle) = \operatorname{vec}_{\Phi}(p)$ for any walk $p$ in $\Gamma$. 
\end{enumerate}
\end{defi}

\setcounter{cdrthm}{6}
\begin{lem}\label{lem:mu}
Let $(\Gamma,L)$ be an $n$-dimensional periodic graph. 
Let $\mu _{\Phi}$ be the homomorphism defined in Definition~\ref{defi:pr}(2). 
The restriction map $\mu_{\Phi} |_{H_1(\Gamma/L, \mathbb{Z})}: H_1(\Gamma/L, \mathbb{Z}) \to L_{\mathbb{R}}$ 
is independent of the choice of the periodic realization $\Phi$. 
Furthermore, its image is contained in $L$. 
\end{lem}
\begin{proof}
Since $H_1(\Gamma/L, \mathbb{Z})$ is generated by $\langle q \rangle$ for $q \in \operatorname{Cyc}_{\Gamma /L}$, 
it is sufficient to show that $\mu_{\Phi} (\langle q \rangle) \in L$ 
and that the value $\mu_{\Phi} (\langle q \rangle)$ is independent of the choice of $\Phi$. 
Take any walk $p$ in $\Gamma$ such that $\overline{p} = q$. 
Since $\overline{t(p)} = \overline{s(p)}$, we have 
\[
\mu_{\Phi} (\langle q \rangle) = \operatorname{vec}_{\Phi}(p) = \Phi(t(p)) - \Phi(s(p)) = t(p) - s(p) \in L, 
\]
and this is independent of the choice of $\Phi$. 
\end{proof}

\begin{defi}\label{defi:mu}
Let $(\Gamma, L)$ be an $n$-dimensional periodic graph. 
We denote by $\mu \colon H_1(\Gamma/L, \mathbb{Z}) \to L$ the restriction of 
$\mu _{\Phi}$ in Definition~\ref{defi:pr}(2) to $H_1(\Gamma/L, \mathbb{Z})$. 
Note that this restriction map is well-defined by Lemma~\ref{lem:mu}. 
\end{defi}

\begin{rmk}\label{rmk:mu}
The homomorphism $\mu$ in Definition~\ref{defi:mu} coincides with $\mu$ defined in~\cite[Section 6.1]{Sunada}. 
\end{rmk}

\setcounter{cdrthm}{10}
\begin{defi}
Let $(\Gamma, L)$ be an $n$-dimensional periodic graph. 
Let $q_0$ be a path in $\Gamma /L$, and let $q_1, \ldots, q_{\ell} \in \operatorname{Cyc}_{\Gamma /L}$ be cycles. 
The sequence $(q_0, q_1, \ldots, q_{\ell})$ is called \textit{walkable} 
if there exists a walk $q'$ in $\Gamma / L$ such that 
$\langle q' \rangle = \sum _{i = 0} ^{\ell} \langle q_i \rangle$. 
\end{defi}

\begin{lem}\label{lem:bunkai}
Let $(\Gamma, L)$ be an $n$-dimensional periodic graph. 
\begin{enumerate}
\item 
For a walk $q'$ in $\Gamma / L$, there exists a walkable sequence $(q_0, q_1, \ldots , q_{\ell})$ such that 
$\langle q' \rangle = \sum _{i = 0} ^{\ell} \langle q_i \rangle$. 

\item 
Let $q_0$ be a path in $\Gamma /L$, and let $q_1, \ldots, q_{\ell} \in \operatorname{Cyc}_{\Gamma /L}$ be cycles. 
Then, $(q_0, q_1, \ldots, q_{\ell})$ is walkable if and only if 
there exists a permutation $\sigma:  \{ 1,2, \ldots, \ell \} \to \{ 1,2, \ldots, \ell \}$ such that
\[
\biggl( \operatorname{supp}(q_0) \cup \bigcup _{1 \le i \le k} \operatorname{supp}(q_{\sigma(i)}) \biggr) 
\cap \operatorname{supp}(q_{\sigma(k+1)}) \not = \emptyset
\] 
holds for any $0 \le k \le \ell -1$.

\end{enumerate}
\end{lem}
\begin{proof}
For any walk $q'$ in $\Gamma /L$, 
if $q'$ is not a path, then 
there exist a walk $q''$ and a cycle $q_1$ in $\Gamma /L$ 
such that $\langle q' \rangle = \langle q'' \rangle + \langle q_1 \rangle$. 
Therefore, the assertion (1) follows by the induction on the length of $q'$. 
The assertion (2) also follows from the induction. 
\end{proof}

\medskip \medskip

\begin{rmk}\label{rmk:bunkai} \hfill
\begin{enumerate}
\item
If $q'$ in Lemma~\ref{lem:bunkai}(1) satisfies $s(q') = t(q')$, then $q_0$ must be a trivial path 
(i.e., $\operatorname{length}(q_0)=0$). 

\item
If a walk $q'$ in $\Gamma /L$ and a vertex $x_0 \in V_{\Gamma}$ satisfy $s(q') = \overline{x_0}$, 
then there exists the unique walk $p$ in $\Gamma$ satisfying $\overline{p} = q'$ and $s(p) = x_0$ 
(we call such $p$ the \textit{lift} of $q'$ with initial point $x_0$). 
Therefore, for a walkable sequence $(q_0, q_1, \ldots, q_{\ell})$ and a vertex $x_0 \in V_{\Gamma}$ satisfying $s(q_0) = \overline{x_0}$, 
there exists a walk $p$ in $\Gamma$ such that 
$s(p) = x_0$ and $\langle \overline{p} \rangle = \sum _{i = 0} ^{\ell} \langle q_i \rangle$. 
Conversely, for a walk $p$ in $\Gamma$, applying Lemma~\ref{lem:bunkai}(1) to $\overline{p}$, 
there exists a walkable sequence $(q_0, q_1, \ldots, q_{\ell})$ satisfying 
$\langle \overline{p} \rangle = \sum _{i = 0} ^{\ell} \langle q_i \rangle$. 

\item
For a walk $p$ in $\Gamma$ and a walkable sequence $(q_0, q_1, \ldots, q_{\ell})$ satisfying 
$\langle \overline{p} \rangle = \sum _{i = 0} ^{\ell} \langle q_i \rangle$,  
it follows that 
\[
w(p) = 
\sum _{i=0}^{\ell} w(q_i), \qquad
\operatorname{length}(p) = 
\sum _{i=0}^{\ell} \operatorname{length}(q_i). 
\]
Furthermore, if we fix a periodic realization $\Phi$, we also have 
\[
\operatorname{vec}_{\Phi}(p) = \sum _{i=0}^{\ell} \mu_{\Phi} ( \langle q_i \rangle). 
\]
\end{enumerate}
\end{rmk}

\setcounter{subsection}{3}
\setcounter{cdrthm}{14}
\subsection{Growth sequences of periodic graphs}\label{subsection:CS}
Let $\Gamma$ be a locally finite graph, and let $x_0 \in V_{\Gamma}$. 
For $i \in \mathbb{Z}_{\ge 0}$, we define subsets $B_{\Gamma, x_0, i}, S_{\Gamma, x_0,i} \subset V _{\Gamma}$ by
\[
B_{\Gamma, x_0,i} \coloneqq  \{ y \in V _{\Gamma} \mid d_{\Gamma}(x_0, y) \le i \}, \quad
S_{\Gamma, x_0,i} \coloneqq  \{ y \in V _{\Gamma} \mid d_{\Gamma}(x_0, y) = i \}. 
\]
Let $b_{\Gamma, x_0,i} \coloneqq  \# B_{\Gamma, x_0,i}$ and $s_{\Gamma, x_0,i} \coloneqq  \# S_{\Gamma, x_0,i}$ denote their cardinalities. 
The sequence $(s_{\Gamma, x_0,i})_i$ is called the \textit{growth sequence} of $\Gamma$ with the start point $x_0$. The sequence $(b_{\Gamma, x_0,i})_i$ is called the \textit{cumulative growth sequence}. 

The \textit{growth series} $G_{\Gamma, x_0}(t)$ is the generating function 
\[
G_{\Gamma, x_0}(t) \coloneqq  \sum _{i \ge 0} s_{\Gamma, x_0,i} t^i
\]
of the growth sequence $(s_{\Gamma, x_0,i})_{i}$. 

\begin{rmk}
In crystallography, the growth sequence is called a \textit{coordination sequence} (see~\cite{NSMN21}). 
\end{rmk}

\begin{defi}[{cf.~\cite[Chapter 0]{Sta96}}]\label{defi:qp}  \hfill
\begin{enumerate}
\item \label{item:f}
A function $f: \mathbb{Z} \to \mathbb{C}$ is called a \textit{quasi-polynomial} 
if there exist a positive integer $N$ and polynomials $Q_0, \ldots, Q_{N-1} \in \mathbb{C}[x]$ such that 
$f(n) = Q_i(n)$ holds for all $n \in \mathbb{Z}$ and $i \in \{ 0, \ldots , N-1 \}$ with $n \equiv i \pmod N$. 
The polynomials $Q_0, \ldots, Q_{N-1}$ are called the \textit{constituents} of $f$. 

\item 
A function $g: \mathbb{Z} \to \mathbb{C}$ is called to be of \textit{quasi-polynomial type} if 
there exists a non-negative integer $M \in \mathbb{Z}_{\ge 0}$ and a quasi-polynomial $f$ such that
$g(n) = f(n)$ holds for all $n > M$. 
The positive integer $N$ is called a \textit{quasi-period} of $g$ when $f$ is of the form in (\ref{item:f}). 
Note that the notion of quasi-period is not unique. 
The minimum quasi-period is called the \textit{period} of $g$. 
We say that the function $g$ is a \textit{quasi-polynomial on} $n \ge m$ if 
$g(n) = f(n)$ holds for $n \ge m$. 
\end{enumerate}
\end{defi}

The growth sequences of periodic graphs are known to be of quasi-polynomial type (Theorem~\ref{thm:NSMN}). 
\begin{thm}[{\cite[Theorem 2.2]{NSMN21}}]\label{thm:NSMN}
Let $(\Gamma, L)$ be a periodic graph, and let $x_0 \in V_{\Gamma}$. 
Then, the functions $b: i \mapsto b_{\Gamma, x_0, i}$ and $s: i \mapsto s_{\Gamma, x_0, i}$ are of quasi-polynomial type. 
In particular, its growth series is a rational function. 
\end{thm}

\noindent
In~\cite{NSMN21}, Theorem~\ref{thm:NSMN} is proved for unweighted periodic graphs, 
and the same proof also works for weighted periodic graphs. 

\setcounter{cdrthm}{18}
\begin{defi}\label{defi:P}
Let $(\Gamma , L)$ be an $n$-dimensional periodic graph. 
\begin{enumerate}
\item 
We define the \textit{normalization map} 
$\nu : \operatorname{Cyc}_{\Gamma /L} \to L_{\mathbb{R}} \coloneqq  L \otimes _{\mathbb{Z}} \mathbb{R}$ by 
\[
\nu: \operatorname{Cyc}_{\Gamma /L} \to L_{\mathbb{R}}; \quad p \mapsto \frac{\mu(\langle p \rangle )}{w(p)}. 
\]
We define the \textit{growth polytope}
\[
P_{\Gamma} \coloneqq  \operatorname{conv} \bigl( \operatorname{Im}(\nu) \cup \{ 0 \} \bigr) \subset L_{\mathbb{R}}
\]
as the convex hull of the set 
$\operatorname{Im}(\nu) \cup \{ 0 \} \subset L_{\mathbb{R}}$. 
Note that $\operatorname{Cyc}_{\Gamma /L}$ is a finite set. 
Furthermore, we have $\operatorname{Im}(\nu) \subset L_{\mathbb{Q}}$ since $w(p) \in \mathbb{Z}_{>0}$ and $\mu(\langle p \rangle ) \in L$ (cf.\ Definition~\ref{defi:mu}). 
Therefore, $P_{\Gamma}$ is a rational polytope 
(i.e., $P_{\Gamma}$ is a polytope whose vertices are on $L_{\mathbb{Q}} \coloneqq  L \otimes _{\mathbb{Z}} \mathbb{Q}$). 
When $\Gamma$ is strongly connected, we have $0 \in \operatorname{int}(P_{\Gamma})$ by Lemma~\ref{lem:int}. 
 
\item 
For a polytope $Q \subset L_{\mathbb{R}}$ and $y \in L_{\mathbb{R}}$, we define 
\[
d_{Q}(y) \coloneqq  \inf \{ t \in \mathbb{R}_{\ge 0} \mid y \in t Q \} \in \mathbb{R}_{\ge 0} \cup \{ \infty \}. 
\]
When $0 \in \operatorname{int}(Q)$, 
we have $d_{Q}(y) < \infty$ for any $y \in L_{\mathbb{R}}$. 

\item
For a periodic realization $\Phi\colon V_{\Gamma} \to L_{\mathbb{R}}$, we define 
\[
d_{P_{\Gamma}, \Phi}(x, y) \coloneqq  d_{P_{\Gamma}}\bigl( \Phi(y) - \Phi(x) \bigr)
\]
for $x, y \in V_{\Gamma}$. 
\end{enumerate}
\end{defi}

\setcounter{cdrthm}{23}
We define $C_{1}(\Gamma, \Phi, x_0)$ and $C_{2}(\Gamma, \Phi, x_0)$ as invariants that measure the difference between $d_{\Gamma}$ and $d_{P_{\Gamma}, \Phi}$. 

\begin{defi}
Let $(\Gamma , L)$ be a strongly connected periodic graph. 
Let $\Phi : V_{\Gamma} \to L_{\mathbb{R}}$ be a periodic realization, and let $x_0 \in V_{\Gamma}$. 
Then, we define 
\begin{align*}
C_{1}(\Gamma, \Phi, x_0) &\coloneqq  \sup _{y \in V_{\Gamma}} \bigl( d_{P_{\Gamma}, \Phi} (x_0, y) - d_{\Gamma} (x_0, y) \bigr), \\
C_{2}(\Gamma, \Phi, x_0) &\coloneqq  \sup _{y \in V_{\Gamma}} \bigl( d_{\Gamma}(x_0, y) - d_{P_{\Gamma}, \Phi} ( x_0, y ) \bigr). 
\end{align*}
By Theorem~\ref{thm:asymp}, we have $C_{1}(\Gamma, \Phi, x_0) < \infty$ and $C_{2}(\Gamma, \Phi, x_0) < \infty$. 
\end{defi}

\section{Ehrhart theory on periodic graphs}\label{section:ETP}

\noindent\textit{Original Section 3 definitions used by the counted claim.}
\begin{defi}
Let $(\Gamma , L)$ be a strongly connected periodic graph, 
and let $\Phi : V_{\Gamma} \to L_{\mathbb{R}}$ be a periodic realization. 
Let $x_0 \in V_{\Gamma}$ and let $\alpha \in \mathbb{R}$. 
The triple $(\Gamma, \Phi, x_0)$ is called to be \textit{$\alpha$-Ehrhart} if we have
\[
B_{\Gamma, x_0, i} = 
\{ y \in V_{\Gamma} \mid d_{P_{\Gamma}, \Phi} (x_0, y) \le i + \alpha \}
\]
for all $i \in \mathbb{Z}_{\ge 0}$. 

This condition is equivalent to the condition that 
\[
d_{P_{\Gamma}, \Phi} (x_0, y) - \alpha \le d_{\Gamma} (x_0, y)
< d_{P_{\Gamma}, \Phi} (x_0, y) + 1 - \alpha
\]
holds for all $y \in V_{\Gamma}$. 
\end{defi}

\begin{defi}
Let $(\Gamma , L)$ be a strongly connected $n$-dimensional periodic graph. 
Let~$\Phi : V_{\Gamma} \to L_{\mathbb{R}}$ be a periodic realization, and let $x_0 \in V_{\Gamma}$. 
We say that \textit{$\Phi$ is symmetric with respect to $x_0$} if 
$\#(\Phi ^{-1}(y)) = \#(\Phi ^{-1}(y'))$ holds for all $y, y' \in L_{\mathbb{R}}$ satisfying $y' + y = 2 \Phi(x_0)$. 
\end{defi}

\begin{rmk}
When $\# (V_{\Gamma} /L) = 1$, there exists an essentially unique periodic realization $\Phi$ (unique up to translation), 
and this $\Phi$ is symmetric with respect to any vertex $x_0 \in V_{\Gamma}$. 
\end{rmk}

\section*{Counted claim: original Theorem 3.4(1)(2) with the complete author proof}
\begin{thm}\label{thm:ET}
Let $(\Gamma , L)$ be a strongly connected $n$-dimensional periodic graph. 
Let~$\Phi\colon V_{\Gamma} \to L_{\mathbb{R}}$ be a periodic realization, and let $x_0 \in V_{\Gamma}$. 
Let $s_i \coloneqq  s_{\Gamma, x_0, i}$ and $b_i \coloneqq  b_{\Gamma, x_0, i}$ be 
the growth sequence and the cumulative growth sequence with the start point $x_0$. 
Let $G_s(t) \coloneqq  \sum _{i \ge 0} s_i t^i$ and $G_b(t) \coloneqq  \sum _{i \ge 0} b_i t^i$ be their generating functions. 
Set $C_1 \coloneqq  C_{1}(\Gamma, \Phi, x_0)$ and $C_2 \coloneqq  C_{2}(\Gamma, \Phi, x_0)$. 
\begin{enumerate}
\item 
Suppose that the triple $(\Gamma, \Phi, x_0)$ is $\alpha$-Ehrhart for some $\alpha \in [0,1)$. 
Then, the function $i \mapsto b_i$ is a quasi-polynomial on $i \ge 0$. 

\item 
Suppose $C_1 + C_2 < 1$. 
Then, $(\Gamma, \Phi, x_0)$ is $\alpha$-Ehrhart for any $\alpha \in [C_{1}, 1-C_{2})$. 

\item 
Suppose both $C_1 < \frac{1}{2}$ and $C_2 < \frac{1}{2}$. 
Suppose one of the following conditions holds: 
\begin{itemize}
\item[(i)] $\Gamma$ is undirected, or

\item[(ii)] $\Phi$ is symmetric with respect to $x_0$. 
\end{itemize}
Then, we have 
\[
G_b(1/t) = (-1)^{n+1} t G_b(t), \qquad 
G_s(1/t) = (-1)^n G_s(t). 
\]
In particular, we have 
\[
f_b(-i) = (-1)^n f_b(i-1)
\]
for any $i \in \mathbb{Z}$, and 
\[
f_s(-i) = (-1)^{n+1} f_s(i)
\]
for any $i \in \mathbb{Z} \setminus \{ 0 \}$, 
where $f_b$ and $f_s$ are the quasi-polynomials corresponding to the sequences $(b_i)_i$ and $(s_i)_i$. 
\end{enumerate}
\end{thm}

\begin{proof}
As in Appendix~\ref{section:ET}, 
for a rational polytope $P \subset L_{\mathbb{R}}$, $v \in L_{\mathbb{R}}$, and $\beta \in \mathbb{R}$, 
we define a function $h_{P, v, \beta}: \mathbb{Z} \to \mathbb{Z}$ by
\[
h_{P,v,\beta}(i) \coloneqq  \# \bigl( (v+(i+\beta)P) \cap L \bigr), 
\]
and its generating function $H_{P,v,\beta}(t) = \sum _{i \in \mathbb{Z}} h_{P,v,\beta}(i)t^i$. 
We also define $\overset{\circ}{h}_{P,v,\beta}$ and $\overset{\circ}{H}_{P,v,\beta}$ similarly. 
Note that we have $\operatorname{relint}(P_{\Gamma}) = \operatorname{int}(P_{\Gamma})$ by the assumption that $\Gamma$ is strongly connected (see Lemma~\ref{lem:int}).


Let $c \coloneqq  \#(V_{\Gamma}/L)$. Take $y_1, \ldots, y_c \in V_{\Gamma}$ such that 
$\{ \overline{y}_1, \ldots , \overline{y}_c \} = V_{\Gamma} /L$. 
Then, we have $V_{\Gamma} = \bigsqcup _{j = 1} ^c (y_j + L)$, and hence, 
\begin{align*}
B_{\Gamma, x_0,i}
&= \left \{ y \in V_{\Gamma} \ \middle | \  d_{P_{\Gamma}, \Phi} (x_0, y) \le i + \alpha \right \} \\
&= \Phi^{-1}\bigl( \Phi (x_0) + (i + \alpha)P_{\Gamma} \bigr) \\
&= \bigsqcup _{j = 1} ^c \left( \Phi^{-1}\bigl( \Phi (x_0) + (i + \alpha)P_{\Gamma} \bigr) \cap (y_j + L) \right) \\
&= \bigsqcup _{j = 1} ^c \bigl( y_j + \bigl \{ m \in L \ \big | \ \Phi(y_j) + m \in \Phi(x_0) +  (i + \alpha)P_{\Gamma} \bigr \} \bigr) \\
&= \bigsqcup _{j = 1} ^c \bigl( y_j + \bigl( \Phi (x_0) - \Phi(y_j) + (i + \alpha) P_{\Gamma} \bigr) \cap L \bigr). 
\end{align*}
Hence, we have 
\[
b_i = \sum _{j=1}^c \# \bigl( \bigl( \Phi (x_0) - \Phi(y_j) + (i + \alpha) P_{\Gamma} \bigr) \cap L \bigr) 
= \sum _{j = 1} ^c h_{P_{\Gamma}, \Phi (x_0) - \Phi(y_j), \alpha}(i). 
\]
Therefore, (1) follows from Theorem~\ref{thm:Eh}(1). 

(2) follows from the definitions of $C_1$ and $C_2$. 

\end{proof}

\noindent\textit{Excluded material (recorded, not claimed).} Assertion (3) of the printed Theorem 3.4, the reciprocity laws for $G_b$ and $G_s$ under $C_1<1/2$ and $C_2<1/2$, is displayed verbatim above only because it belongs to the printed statement; it is not claimed here, and its published proof (source lines 1072--1128) is not part of this excerpt. Remarks 3.5--3.12, the polytope-to-graph construction, and all later examples, tables and figures are likewise excluded and are not used to inflate the difficulty judgement.
\appendix
\section{Properties of the growth polytope}\label{section:GP}

\noindent\textit{Uncounted context and required core.} Lemma A.1 is quoted with its complete published proof because the counted proof uses the identity $\operatorname{relint}(P_\Gamma)=\operatorname{int}(P_\Gamma)$ for strongly connected periodic graphs. The statement of Theorem A.2 is quoted because Definition 2.24 invokes it for the finiteness of $C_1$ and $C_2$; its published proof is not part of this excerpt. Appendix B is the complete required core for the counted claim.
\begin{lem}[{\cite[Proposition 21]{Fri13}}]\label{lem:int}
If $(\Gamma, L)$ is a strongly connected $n$-dimensional periodic graph, then 
we have $0 \in \operatorname{int}(P _{\Gamma})$. 
In particular, we have $d_{P _{\Gamma}}(y) < \infty$ for any $y \in L_{\mathbb{R}}$. 
\end{lem}
\begin{proof}
It is sufficient to prove the following claim: 
\begin{itemize}
\item For any $u \in L$, there exists $m \in \mathbb{Z}_{>0}$ such that $u \in m P _{\Gamma}$. 
\end{itemize}
We fix $x_0 \in V_{\Gamma}$. 
Since the graph $\Gamma$ is strongly connected by assumption, 
there exists a walk $p$ in $\Gamma$ from $x_0$ to $u + x_0$. 
Then, by applying Lemma~\ref{lem:bunkai}(1) to $\overline{p}$ (cf.\ Remark~\ref{rmk:bunkai}(1)), 
there exists a sequence $q_1, \ldots, q_{\ell} \in \operatorname{Cyc}_{\Gamma /L}$ satisfying
$\langle \overline{p} \rangle = \sum _{i=1} ^{\ell} \langle q_i \rangle$. 
Since $\frac{\mu (\langle q \rangle)}{w(q)} \in P _{\Gamma}$ holds for any $q \in \operatorname{Cyc}_{\Gamma /L}$, 
we have 
\[
u = \mu ( \langle \overline{p} \rangle ) = \sum _{i=1}^{\ell} \mu(\langle q_i \rangle)
\in \biggl( \sum _{i=1} ^{\ell} w(q_i) \biggr) P _{\Gamma}, 
\]
which completes the proof. 
\end{proof}

\setcounter{cdrthm}{1}
\begin{thm}\label{thm:asymp}
Let $(\Gamma , L)$ be a strongly connected $n$-dimensional periodic graph. 
Let $\Phi\colon V_{\Gamma} \to L_{\mathbb{R}}$ be a periodic realization, 
and let $x_0 \in V_{\Gamma}$. Then there exist $C' _1, C' _2 \in \mathbb{R}_{\ge 0}$ such that 
for any $y \in V _{\Gamma}$, we have 
\[
d_{P_{\Gamma}, \Phi} (x_0, y) - C' _1 \le d_{\Gamma}(x_0, y) \le d_{P_{\Gamma}, \Phi} (x_0,y) + C' _2. 
\]
\end{thm}

\section{Ehrhart theory}\label{section:ET}
In this section, we discuss a variant of Ehrhart theory (Theorem~\ref{thm:Eh}), which is necessary for the proof of Theorem~\ref{thm:ET}. The difference from the usual Ehrhart theory is that the center $v$ of the dilation need not be the origin, and the dilation factor may be shifted by a constant $\alpha$. 

\begin{defi}
Let $P \subset \mathbb{R}^N$ be a rational polytope, and let $v \in \mathbb{R}^N$ and $\alpha \in \mathbb{R}$. 

\begin{enumerate}
\item
We define a function $h_{P, v, \alpha}: \mathbb{Z} \to \mathbb{Z}$ as follows
\[
h_{P, v, \alpha}(d) \coloneqq  \# \left( (v + (d + \alpha) P) \cap \mathbb{Z}^N \right). 
\]
Here, we define $t P = \emptyset$ for $t < 0$. 

\item
Let $\operatorname{relint}(P)$ denote the relative interior of $P$. 
Then, we also define a function $\overset{\circ}{h} _{P, v, \alpha}: \mathbb{Z} \to \mathbb{Z}$ as follows
\[
\overset{\circ}{h}_{P, v, \alpha}(d) \coloneqq  \# \left( (v + (d + \alpha) \cdot \operatorname{relint}({P})) \cap \mathbb{Z}^N \right). 
\]
Here, we define $t \cdot \operatorname{relint}({P}) = \emptyset$ for $t \le 0$. 

\item
Let  
\[
H_{P, v, \alpha} (t) \coloneqq  \sum _{i \in \mathbb{Z}} h_{P, v, \alpha}(i) t^i, \qquad
\overset{\circ}{H}_{P, v, \alpha} (t) \coloneqq  \sum _{i \in \mathbb{Z}} \overset{\circ}{h}_{P, v, \alpha}(i) t^i
\]
denote their generating functions. 
\end{enumerate}
\end{defi}

\begin{rmk}  \hfill
\begin{enumerate}
\item
In the definition above, we have defined $0 \cdot P = \{ 0 \}$ but $0 \cdot \operatorname{relint}({P}) = \emptyset$. This definition is necessary for the equations ($\spadesuit$) in the proof of Theorem~\ref{thm:Eh}.

\item
The usual Ehrhart theory (cf.~\cite{BR}) treats the case where $v = 0$ and $\alpha = 0$. 
McMullen (in~\cite{McM78}) and de Vries and Yoshinaga (in~\cite[Section 3]{dVY21}) discuss $h_{P, v, \alpha}$ when $\alpha = 0$. 
\end{enumerate}
\end{rmk}

\begin{lem}\label{lem:sym}
We have $h_{P,v,\alpha} = h_{-P,-v,\alpha}$ and $\overset{\circ}{h}_{P,v,\alpha} = \overset{\circ}{h}_{-P,-v,\alpha}$. 
\end{lem}
\begin{proof}
The first assertion follows from
\begin{align*}
\# \left( (v + (d + \alpha) P) \cap \mathbb{Z}^N \right) 
&= \# \left( (- (v + (d + \alpha) P)) \cap \mathbb{Z}^N \right) \\
&= \# \left( (-v + (d + \alpha) (-P) \cap \mathbb{Z}^N \right). 
\end{align*}
The second assertion can be proved in the same way.
\end{proof}

\begin{thm}\label{thm:Eh}
Let $P \subset \mathbb{R}^N$ be a rational polytope of dimension $M$, and let $v \in \mathbb{R}^N$ and $\alpha \in \mathbb{R}$. 
Then, the following assertions hold: 
\begin{enumerate}
\item $h_{P, v, \alpha}$ is a quasi-polynomial on $d \ge - \alpha$, and 
$h_{P, v, \alpha}(d) = 0$ holds for $d < - \alpha$. 

\item $\overset{\circ}{h} _{P, v, \alpha}$ is a quasi-polynomial on $d > - \alpha$, and 
$\overset{\circ}{h} _{P, v, \alpha}(d) = 0$ holds for $d \le - \alpha$. 

\item $H_{P, v, \alpha}(1/t) = (-1)^{M+1} \overset{\circ}{H}_{P, -v, - \alpha}(t)$. 

\item 
Let $f_{P,v,\alpha}$ be the corresponding quasi-polynomial to the function $h_{P, v, \alpha}$ on $d \ge - \alpha$. 
Then, we have 
$f_{P,v,\alpha}(-i) = (-1)^M \overset{\circ}{h}_{P, -v, - \alpha}(i)$ for any $i \in \mathbb{Z}_{> \alpha}$. 
\end{enumerate}
\end{thm}

\begin{proof}
For a subset $S \subset \mathbb{R}^{N+1}$, let 
\[
\sigma _S({\bf z}) = \sigma _S (z_1, \ldots, z_{N+1}) \coloneqq  \sum _{{\bf m} \in S \cap \mathbb{Z}^{N+1}} {\bf z}^{\bf m}
\]
denote the integer-point transform of $S$ (see~\cite[Section 3.3]{BR}). 
This is a formal sum of Laurent monomials ${\bf z}^{\bf m} = z_1^{m_1} \cdots z_{N+1}^{m_{N+1}}$ 
for ${\bf m} = (m_1, \ldots , m_{N+1}) \in S \cap \mathbb{Z}^{N+1}$. 

Let $K \coloneqq  \mathbb{R}_{\ge 0} \bigl( \{ 1 \} \times P \bigr) \subset \mathbb{R}^{N+1}$ be the cone over $P$. 
Then, we have 
\begin{align*}
\begin{split}
H_{P,v,\alpha} (t) &= \sigma _{(- \alpha, v) + K} (t, 1 , \ldots, 1), \\
\overset{\circ}{H}_{P, -v, - \alpha}(t) &= \sigma _{(\alpha, - v) + \operatorname{relint}(K)} (t, 1 , \ldots, 1). 
\end{split} \tag{$\spadesuit$}
\end{align*}
Furthermore, we have 
\[
\sigma _{(- \alpha, v) + K} (z_1^{-1}, \ldots, z_{N+1}^{-1}) 
= (-1)^{M+1} \sigma _{(\alpha, -v) + \operatorname{relint}(K)} (z_1, \ldots, z_{N+1})
\]
by~\cite[Exercises 4.5 and 4.6]{BR}. Therefore, we get the desired equality in (3). 

The second assertions of (1) and (2) are obvious from the definition of $h_{P, v, \alpha}$ and $\overset{\circ}{h}_{P, v, \alpha}$. 

By taking a triangulation of $P$, 
the first assertions in (1) and (2) can be reduced to the same assertions for simplices $P$. 
Assume that $P$ is a simplex. Let $v_1, \ldots, v_{M+1}$ be the vertices of $P$. 
Since $P$ is rational polytope, we may write $v_i = u_i / a_i$ for some $u_i \in \mathbb{Z}^N$ and $a_i \in \mathbb{Z}_{>0}$. 
We set $D, D^{\circ} \subset \mathbb{R}^{N+1}$ by
\begin{align*}
D &\coloneqq  \left\{ \sum _{i = 1}^{M+1}  \alpha _i (a_i, u_i) \ \middle | \ \alpha_1, \ldots, \alpha_{M+1} \in [0,1) \right\}, \\ 
D^{\circ} &\coloneqq  \left\{ \sum _{i = 1}^{M+1} \alpha _i (a_i, u_i) \ \middle | \ \alpha_1, \ldots, \alpha_{M+1} \in (0,1] \right\}.
\end{align*}
Then, by~\cite[Theorem 3.5]{BR}, we have 
\begin{align*}
\sigma _{(- \alpha, v) + K} (z_1, {\bf z}) 
&= \frac{\sigma _{(- \alpha, v) + D} (z_1, {\bf z})}{\prod _{i=1} ^{M+1} (1-z_1^{a_i} {\bf z}^{u_i})}, \\
\sigma _{(- \alpha, v) + \operatorname{relint}(K)} (z_1, {\bf z}) 
&= \frac{\sigma _{(- \alpha, v) + D^{\circ}} (z_1, {\bf z})}{\prod _{i=1} ^{M+1} (1-z_1^{a_i} {\bf z}^{u_i})}, 
\end{align*}
where ${\bf z} = (z_2, \ldots, z_{N+1})$. 
Therefore, we have 
\begin{align*}
H_{P,v,\alpha} (t) 
&= \sigma _{(- \alpha, v) + K} (t, 1, \ldots , 1) 
= \frac{\sigma _{(- \alpha, v) + D} (t, 1, \ldots, 1)}{\prod _{i=1} ^{M+1} (1- t^{a_i})}, \\
\overset{\circ}{H}_{P,v,\alpha} (t) 
&= \sigma _{(- \alpha, v) + \operatorname{relint}(K)} (t, 1, \ldots , 1) 
= \frac{\sigma _{(- \alpha, v) + D^{\circ}} (t, 1, \ldots, 1)}{\prod _{i=1} ^{M+1} (1- t^{a_i})}. 
\end{align*}
Here, we have 
\begin{align*}
\deg \sigma _{(- \alpha, v) + D} (t, 1, \ldots, 1) 
&< - \alpha + \sum _{i=1} ^{M+1} a_i, \\
\deg \sigma _{(- \alpha, v) + D^{\circ}} (t, 1, \ldots, 1) 
&\le - \alpha + \sum _{i=1} ^{M+1} a_i. 
\end{align*}
Therefore, we can conclude that $h_{P, v, \alpha}$ is a quasi-polynomial on $d \ge - \alpha$, and 
that $\overset{\circ}{h}_{P, v, \alpha}$ is a quasi-polynomial on $d > - \alpha$. 
We complete the proof of (1) and (2). 

Finally, we prove (4). 
Since $f_{P,v,\alpha}$ is a quasi-polynomial, we have
\[
\sum _{i \in \mathbb{Z}_{< - \alpha}} {f_{P,v,\alpha}(i)} t^{i} 
= - \sum _{i \in \mathbb{Z}_{\ge - \alpha}} {f_{P,v,\alpha}(i)} t^i
\]
as rational functions (cf.~\cite[Exercise 4.7]{BR}). 
Therefore, we have 
\begin{align*}
\sum _{i \in \mathbb{Z}_{> \alpha}} {f_{P,v,\alpha}(- i)} t^{- i}
&= \sum _{i \in \mathbb{Z}_{< - \alpha}} {f_{P,v,\alpha}(i)} t^{i} \\ 
&=- \sum _{i \in \mathbb{Z}_{\ge - \alpha}} {f_{P,v,\alpha}(i)} t^i \\
&=- H_{P,v,\alpha}(t) \\
&= (-1)^{M} \overset{\circ}{H}_{P, -v, - \alpha}(t^{-1}) \\
&= (-1)^{M} \sum _{i \in \mathbb{Z}} \overset{\circ}{h}_{P, -v, - \alpha}(i) t^{-i}. 
\end{align*}
Here, the third equality follows from (1), and the fourth follows from (3). 
By comparing the coefficients, we conclude that $f_{P,v,\alpha}(- i) = (-1)^M \overset{\circ}{h}_{P, -v, - \alpha}(i)$ for $i \in \mathbb{Z}_{> \alpha}$. 
\end{proof}

\begin{thebibliography}{99}
\bibitem[5]{BR} Matthias Beck and Sinai Robins, \emph{Computing the continuous discretely: integer-point enumeration in polyhedra}, second ed., Undergraduate Texts in Mathematics, Springer, New York, 2015, https://doi.org/10.1007/978-1-4939-2969-6.
\bibitem[9]{dVY21} Christopher de Vries and Masahiko Yoshinaga, ``Ehrhart quasi-polynomials of almost integral polytopes'', forthcoming, Discrete Comput. Geom., 2021, https://arxiv.org/abs/2108.11132.
\bibitem[11]{Fri13} Tobias Fritz, ``Velocity polytopes of periodic graphs and a no-go theorem for digital physics'', \emph{Discrete Math.} \textbf{313} (2013), no. 12, 1289--1301, https://doi.org/10.1016/j.disc.2013.02.010.
\bibitem[17]{McM78} P. McMullen, ``Lattice invariant valuations on rational polytopes'', \emph{Arch. Math. (Basel)} \textbf{31} (1978/79), no. 5, 509--516, https://doi.org/10.1007/BF01226481.
\bibitem[19]{NSMN21} Yusuke Nakamura, Ryotaro Sakamoto, Takafumi Mase, and Junichi Nakagawa, ``Coordination sequences of crystals are of quasi-polynomial type'', \emph{Acta Crystallogr. A} \textbf{77} (2021), no. 2, 138--148, https://doi.org/10.1107/S2053273320016769.
\bibitem[24]{Sta96} Richard P. Stanley, \emph{Combinatorics and commutative algebra}, second ed., Progress in Mathematics, vol. 41, Birkh\"auser Boston, Inc., Boston, MA, 1996.
\bibitem[25]{Sunada} Toshikazu Sunada, \emph{Topological crystallography: with a view towards discrete geometric analysis}, Surveys and Tutorials in the Applied Mathematical Sciences, vol. 6, Springer, Tokyo, 2013, https://doi.org/10.1007/978-4-431-54177-6.
\end{thebibliography}
\end{document}
